\chapter{Teknikker til kodeanalyse}\label{chap:ch15}

I softwareudvikling er det afgørende at forstå kodens effektivitet, hvis man vil bygge skalerbare applikationer med hurtige svartider. Med udgangspunkt i principperne for asymptotisk notation går vi nu over til praktisk kodeanalyse -- vi skiller koden ad og ser på, hvordan komponenter som løkker, betingelser og sekventielle sætninger bidrager til en algoritmes samlede kompleksitet.

Kodeanalyse kan groft inddeles i to kategorier:

\begin{enumerate}
    \item \textbf{Statisk analyse:} Her undersøger vi koden uden at køre den. Vi fokuserer på at vurdere kodens struktur og operationer, især gennem kompleksitetsanalyse. Det er dette kapitels hovedemne.
    \item \textbf{Dynamisk analyse:} Ved at køre koden på forskellige input måler vi de faktiske køretider og det faktiske hukommelsesforbrug i bestemte miljøer. Dynamisk analyse er værdifuld, men den ligger uden for dette kapitels rammer og bruges typisk til ydelsesprofilering.
\end{enumerate}

Asymptotisk notation giver et overordnet billede af en algoritmes ydelse, idet den forudsiger algoritmens vækstrate, når inputstørrelsen vokser. Den teoretiske kompleksitet tager dog ikke højde for konstanter, mindre led eller maskinspecifikke faktorer. I praksis kan disse detaljer have stor betydning for køretiden, især i applikationer, der ofte udfører operationer på store datasæt.

Se på to algoritmer, der begge har køretiden $\Theta(n)$. De har samme asymptotiske vækstrate, men den ene algoritme kan have en konstant faktor, der er flere gange større end den andens, så den i praksis kører langsommere. For at forstå kompleksitet i praksis må vi analysere hver del af koden, vurdere antallet af operationer konkret og finde de steder, hvor koden kan optimeres.

Nedenfor ser du to algoritmer, der filtrerer de lige tal ud af et heltalsarray. Begge algoritmer har køretiden $\Theta(n)$, men den ene bruger en enkelt løkke, mens den anden foretager et overflødigt indledende gennemløb og derfor har en større konstant faktor i denne model.


\begin{lstlisting}[style=javaStyle, caption={Ét gennemløb til at filtrere lige tal}, label={lst:java_single_pass1}]
    public List<Integer> filterEvenSinglePass(int[] arr) {
        List<Integer> evenNumbers = new ArrayList<>();
        for (int i = 0; i < arr.length; i++) {
            if (arr[i] % 2 == 0) {
                evenNumbers.add(arr[i]);
            }
        }
        return evenNumbers;
    }
\end{lstlisting}

\newpage

\begin{lstlisting}[style=javaStyle, caption={Flere gennemløb til at filtrere lige tal}, label={lst:java_multiple_passes}]
    public List<Integer> filterEvenMultiplePasses(int[] arr) {
        List<Integer> evenNumbers = new ArrayList<>();

        // First pass: count even numbers
        int evenCount = 0;
        for (int i = 0; i < arr.length; i++) {
            if (arr[i] % 2 == 0) {
                evenCount++;
            }
        }

        // Second pass: collect even numbers
        for (int i = 0; i < arr.length; i++) {
            if (arr[i] % 2 == 0) {
                evenNumbers.add(arr[i]);
            }
        }

        return evenNumbers;
    }
\end{lstlisting}

Algoritme \ref{lst:java_single_pass1} gennemløber arrayet én gang for at finde de lige tal. Undervejs tilføjer den hvert lige tal direkte til en resultatliste, så der kun er ét gennemløb på $\mathcal{O}(n)$.

Algoritme \ref{lst:java_multiple_passes} foretager derimod to fulde gennemløb af arrayet. Det første gennemløb tæller antallet af lige elementer, selv om dette antal ikke bruges bagefter, og det andet gennemløb samler elementerne. Begge algoritmer har køretiden $\Theta(n)$, men det overflødige første gennemløb giver ekstra arbejde. Forholdet mellem dem er ikke blot to, fordi listeoperationer, sammenligninger og andre sætninger også bidrager til køretiden.

I praksis er Algoritme \ref{lst:java_single_pass1} hurtigere, fordi den løser opgaven i ét gennemløb, mens Algoritme \ref{lst:java_multiple_passes} får en større konstant faktor ved at gennemløbe arrayet to gange. Eksemplet viser, hvordan forskellige implementeringer med samme kompleksitet, $\mathcal{O}(n)$, alligevel kan yde forskelligt i praksis på grund af forskellige konstante faktorer.

I dette kapitel analyserer vi systematisk kodens komponenter. Vi begynder med pseudokode, så vi kan se bort fra sprogspecifik syntaks. Derefter undersøger vi kompleksiteten af forskellige kodestrukturer, herunder løkker, betingelser og sekvenser. Til sidst ser vi på praktiske forhold og diskuterer, hvordan konstanter og led af lavere orden, som man ofte ser bort fra i asymptotisk notation, kan påvirke ydelsen i praksis.


\section{Pseudokode}\label{sec:ch11-1}
For at kunne analysere algoritmer effektivt uden at blive tynget af sprogspecifik syntaks bruger vi pseudokode -- en forenklet, overordnet fremstilling af kode, der fremhæver en algoritmes logik og struktur. Med pseudokode kan vi beskrive og analysere algoritmer uafhængigt af et bestemt programmeringssprog og udelukkende fokusere på rækkefølgen af operationerne og deres effektivitet. Pseudokode fungerer som et bindeled mellem det konceptuelle algoritmedesign og implementeringen. Fordi sprogets detaljer er skåret væk, gør pseudokode det lettere at se en algoritmes centrale operationer og dermed lettere at vurdere kompleksiteten og optimere ydelsen, før algoritmen implementeres i kode.

Pseudokode gør det også lettere at samarbejde, fordi teammedlemmer med forskellige sprogkompetencer kan forstå og bidrage til algoritmedesignet uden at kende en bestemt syntaks. Desuden er pseudokode nyttig til dokumentation af algoritmer, da den giver en kortfattet, men klar opsummering af de enkelte trin, som er lettere at læse end fuldt udfoldet kode.

Vi følger konventionerne fra \cite{clrs4}, som består af et bestemt sæt formateringsregler, der sikrer klarhed og ensartethed:

\begin{enumerate}
    \item \textbf{Løkkestrukturer}: Vi bruger \texttt{\For} og \texttt{\While} til at angive iterationer:
    \begin{itemize}
        \item \texttt{\For $i \gets 1$ \To $n$} betegner en løkke, der øger \( i \) med 1 fra 1 til og med \( n \).
        \item \texttt{\While} betingelse betegner en løkke, der fortsætter, så længe den angivne betingelse er sand.
    \end{itemize}
    \item \textbf{Tildeling og operationer}: Tildelingsoperatoren $\gets$ angiver, at en variabel tildeles en værdi, fx sætter \texttt{$sum \gets 0$ } \texttt{sum} til nul. De grundlæggende regneoperationer (\texttt{+, -, *, /}) skrives med den sædvanlige notation.
    \item \textbf{Kald af procedurer og funktioner}: Procedurenavne skrives med kapitæler, fx \texttt{$\proc{Sum-Array}(arr)$}, mens variabler og identifikatorer skrives med kursiv.
    \item \textbf{Kommentarer}: Kommentarer vises med \textcolor{commentcolor}{bordeauxrød} farve og giver kontekstuel information om hvert trin.
\end{enumerate}

Lad os som illustration se på en simpel algoritme, der beregner summen af elementerne i et array. Her bruger vi pseudokode til at beskrive algoritmen med fokus på logikken og de enkelte trin.

\begin{codebox}
    \Procname{$\proc{Sum-Array}(arr)$}
    \li $sum \gets 0$ \CommentAt{7}{initialise sum to zero}
    \li \For $i \gets 0$ \To $\attrib{arr}{length} - 1$ \Do
    \li     $sum \gets sum + arr[i]$ \CommentAt{7}{Add the current element to sum}
        \End
    \li \Return $sum$ \CommentAt{7}{Return the final sum}
\end{codebox}

I denne pseudokode:
\begin{itemize}
    \item Proceduren \texttt{$\proc{Sum-Array}$} tager et array \texttt{arr} som input.
    \item Vi initialiserer en variabel \texttt{sum} til nul.
    \item En \texttt{\For}-løkke gennemløber hvert element i arrayet og lægger det til \texttt{sum}.
    \item Til sidst returnerer funktionen den beregnede \texttt{sum}.
\end{itemize}

Pseudokoden indfanger algoritmens væsentlige trin uden syntaksspecifikke detaljer, så den er let at analysere og tilpasse. Her er et mere komplekst eksempel:

\begin{codebox}
    \Procname{$\proc{Find-Max-Even-Sum}(arr)$}
    \li $maxSum \gets 0$ \CommentAt{7}{Initialise max sum to zero}
    \li $currentSum \gets 0$ \CommentAt{7}{Track the sum of the current subarray}
    \li \For $i \gets 0$ \To $\attrib{arr}{length} - 1$ \Do
    \li \If $arr[i] \bmod 2 == 0$ \Then
    \li $currentSum \gets currentSum + arr[i]$ \CommentAt{12}{Add even number to current sum}
    \li $maxSum \gets \proc{max}(maxSum, currentSum)$ \CommentAt{12}{Update if currentSum > maxSum}
    \li \Else \li $currentSum \gets 0$ \CommentAt{7}{Reset current sum if odd number is encountered} \End \End
    \li \Return $maxSum$ \CommentAt{7}{Return the maximum sum of any contiguous even subarray}
\end{codebox}

Proceduren beregner den største sum af sammenhængende lige elementer i et array:

\begin{itemize}
    \item Proceduren tager et inputarray \texttt{arr} og initialiserer \texttt{maxSum}, der holder styr på den største sum, og \texttt{currentSum}, der opsamler summen af det aktuelle sammenhængende delarray af lige tal.
    \item En enkelt \texttt{\For}-løkke gennemløber \texttt{arr}:
    \item Hvis et element er lige, lægges det til \texttt{currentSum}, og \texttt{maxSum} opdateres, hvis \texttt{currentSum} overstiger den.
    \item Hvis et element er ulige, nulstilles \texttt{currentSum} til 0, fordi det sammenhængende delarray er afbrudt.
\end{itemize}

Her er yderligere to eksempler på pseudokode, et for \textit{iterativ binær søgning }og et for \textit{rekursiv binær søgning}. Pseudokoden for de to algoritmer ligner hinanden, men den rekursive version kalder sig selv via et funktionskald, mens den iterative version styrer søgningen med en løkke.

\begin{codebox}
    \Procname{$\proc{Binary-Search-Iterative}(arr, key)$}
    \li $low \gets 0$ \CommentAt{7}{Initialise low index}
    \li $high \gets arr.length -1$ \CommentAt{7}{Initialise high index}
    \li \While $low \leq high$ \Do
    \li     $mid \gets \floor{(low + high) / 2}$ \CommentAt{7}{Calculate mid index}
    \li     \If $arr[mid] == key$ \Then
    \li         \Return $mid$ \CommentAt{7}{Key found at mid index}
    \li     \ElseIf $arr[mid] < key$ \Then
    \li         $low \gets mid + 1$ \CommentAt{7}{Update low index}
    \li     \Else
    \li         $high \gets mid - 1$ \CommentAt{7}{Update high index}
        \End
        \End
    \li \Return $-1$ \CommentAt{7}{Key not found}
\end{codebox}

Her er en trinvis forklaring af, hvad koden gør:

\begin{enumerate}
    \item \textbf{Initialisering}:
        \begin{itemize}
            \item \texttt{low} initialiseres til det første indeks i arrayet \texttt{arr}.
            \item \texttt{high} initialiseres til det sidste indeks i arrayet \texttt{arr}.
        \end{itemize}
    \item \textbf{While-løkke}:
        \begin{itemize}
            \item Løkken fortsætter, så længe \texttt{low} er mindre end eller lig med \texttt{high}.
        \end{itemize}
    \item \textbf{Beregning af midterindeks}:
        \begin{itemize}
            \item \texttt{mid} beregnes som gennemsnittet af \texttt{low} og \texttt{high} rundet ned.
        \end{itemize}
    \item \textbf{Sammenligning}:
        \begin{itemize}
            \item Hvis elementet på indeks \texttt{mid} er lig med \texttt{key}, returnerer funktionen \texttt{mid}, som angiver positionen af \texttt{key} i arrayet.
            \item Hvis elementet på indeks \texttt{mid} er mindre end \texttt{key}, må \texttt{key} ligge i arrayets højre halvdel. Derfor opdateres \texttt{low} til \texttt{mid + 1}.
            \item Hvis elementet på indeks \texttt{mid} er større end \texttt{key}, må \texttt{key} ligge i arrayets venstre halvdel. Derfor opdateres \texttt{high} til \texttt{mid - 1}.
        \end{itemize}
    \item \textbf{Nøglen findes ikke}:
        \begin{itemize}
            \item Hvis løkken afsluttes uden at have fundet \texttt{key}, returnerer funktionen \texttt{-1}, hvilket angiver, at \texttt{key} ikke findes i arrayet.
        \end{itemize}
\end{enumerate}

Algoritmen søger effektivt efter en nøgle, \texttt{key}, i et sorteret array ved gentagne gange at halvere søgeintervallet. Tidskompleksiteten for binær søgning er \(\mathcal{O}(\log n)\), så den er meget hurtigere end lineær søgning for store arrays. Algoritmen er \textbf{iterativ}, fordi den bruger en \texttt{\While}-løkke til at udføre en kodeblok gentagne gange, indtil en bestemt betingelse er opfyldt. Her fortsætter løkken, så længe \texttt{low} er mindre end eller lig med \texttt{high}. I hver iteration af løkken udføres en sammenligning, og indeksene \texttt{low} og \texttt{high} opdateres tilsvarende. Denne iterative fremgangsmåde står i kontrast til en rekursiv fremgangsmåde, hvor funktionen kalder sig selv med opdaterede parametre, indtil basisbetingelsen er opfyldt.

\begin{codebox}
    \Procname{$\proc{Binary-Search-Recursive}(arr, key, low, high)$}
    \li \If $low > high$ \Then
    \li     \Return $-1$ \CommentAt{7}{Key not found} \End
    \li $mid \gets \floor{(low + high) / 2}$ \CommentAt{7}{Calculate mid index}
    \li \If $arr[mid] == key$ \Then
    \li     \Return $mid$ \CommentAt{7}{Key found at mid index}
    \li \ElseIf $arr[mid] < key$ \Then
    \li     \Comment{Search right half}
    \li     \Return $\proc{Binary-Search-Recursive}(arr, key, mid + 1, high)$
    \li \Else
    \li     \Comment{Search left half}
    \li     \Return $\proc{Binary-Search-Recursive}(arr, key, low, mid - 1)$ 
        \End
\end{codebox}

Koden implementerer en rekursiv binær søgealgoritme. Funktionen $\proc{Binary-Search-Recursive}$ tager et array \texttt{arr}, en nøgle \texttt{key}, der skal søges efter, og indeksene \texttt{low} og \texttt{high}, som afgrænser det aktuelle søgeområde.

\begin{enumerate}
    \item \textbf{Basistilfælde}: Hvis \texttt{low} er større end \texttt{high}, returnerer funktionen $-1$, hvilket angiver, at nøglen ikke findes i arrayet.
    \item \textbf{Beregning af midtpunkt}: Midtpunktet \texttt{mid} beregnes som gennemsnittet af \texttt{low} og \texttt{high} rundet ned.
    \item \textbf{Sammenligning med nøglen}:
    \begin{itemize}
        \item Hvis elementet på indeks \texttt{mid} er lig med \texttt{key}, returnerer funktionen \texttt{mid}, hvilket angiver, at nøglen er fundet på dette indeks.
        \item Hvis elementet på indeks \texttt{mid} er mindre end \texttt{key}, søger funktionen rekursivt i arrayets højre halvdel ved at opdatere \texttt{low} til \texttt{mid + 1}.
        \item Hvis elementet på indeks \texttt{mid} er større end \texttt{key}, søger funktionen rekursivt i arrayets venstre halvdel ved at opdatere \texttt{high} til \texttt{mid - 1}.
    \end{itemize}
\end{enumerate}

Algoritmen søger effektivt efter en nøgle, \texttt{key}, i et sorteret array ved gentagne gange at halvere søgeintervallet. Algoritmen er \textbf{rekursiv}, fordi den kalder sig selv med opdaterede parametre for at indsnævre søgeområdet. I hvert rekursivt kald beregner funktionen midtpunktet af det aktuelle område og foretager en sammenligning. Afhængigt af resultatet returnerer funktionen enten indekset \texttt{mid}, hvis nøglen er fundet, eller kalder sig selv rekursivt på enten venstre eller højre halvdel af det aktuelle interval, idet indekset \texttt{low} eller \texttt{high} justeres tilsvarende. Rekursionen fortsætter, indtil basisbetingelsen er opfyldt, dvs. indtil nøglen enten er fundet, eller det er fastslået, at den ikke findes. Denne rekursive fremgangsmåde står i kontrast til en iterativ fremgangsmåde, der udfører søgningen med en løkke i stedet for rekursive funktionskald.

\section{Analyse af kodekonstruktioner}\label{sec:ch11-2}
Ved konstruktionsanalyse deler vi koden op i dens grundlæggende komponenter, fx løkker, betingelser og sekvenser, for at forstå, hvordan hver af dem bidrager til den samlede kompleksitet. Ved at analysere disse konstruktioner kan vi finde flaskehalse, ineffektiviteter og steder i en algoritme, der kan optimeres. I dette afsnit undersøger vi kompleksiteten af forskellige kodestrukturer med fokus på løkker, betingelser og sekvenser, og vi ser på, hvordan de påvirker en algoritmes ydelse.

\subsection*{Basisoperationer}
Den tid, en funktion eller procedure bruger, er proportional med det antal »basisoperationer«, den udfører. Eksempler på basisoperationer er:

\begin{enumerate}
    \item \textbf{Regneoperation}: En enkelt regneoperation, fx addition (+), multiplikation (\texttt{*}), subtraktion (-) eller division (/).
    \item \textbf{Tildeling}: At tildele en variabel en værdi, fx $x \gets 0$.
    \item \textbf{Test/sammenligning}: At evaluere en betingelse, fx $x == 0$ eller $i \leq n$.
    \item \textbf{Funktionskald}: At kalde en funktion eller procedure, fx et kald af \proc{Find-Max-Even-Sum}.
    \item \textbf{Return-sætning}: At returnere en værdi fra en funktion, fx \texttt{return total}.
    \item \textbf{Læseoperation}: At læse en værdi af en primitiv type (fx heltal, kommatal, tegn eller boolesk værdi) fra input eller hukommelsen.
    \item \textbf{Skriveoperation}: At skrive en værdi af en primitiv type til output eller hukommelsen.
\end{enumerate}

Disse basisoperationer er grundlaget for mere komplekse operationer. I den forenklede \emph{enhedsomkostningsmodel}, som vi bruger i dette kapitel, tildeles hver operation på primitive værdier af fast størrelse en konstant omkostning. Antagelsen gælder ikke uændret for vilkårligt store heltal eller for operationer, hvis omkostning afhænger af, hvordan data er repræsenteret. Ved at tælle operationer i denne model kan vi vurdere en konstruktions køretid. I vores detaljerede optællinger medtager vi kun punkt 1--5 og ser bort fra særskilte omkostninger til læsning og skrivning.

Alle de følgende algoritmer er $\mathcal{O}(1)$-algoritmer, fordi de udfører et konstant antal basisoperationer:

\begin{codebox}
    \Procname{$\proc{Inc}(x)$}
    \li \Return $x + 1$
\end{codebox}

\begin{codebox}
    \Procname{$\proc{Mul}(x, y)$}
    \li \Return $x \times y$
\end{codebox}

\begin{codebox}
    \Procname{$\proc{Foo}(x)$}
    \li $y \gets x \times 77.3$
    \li \Return $x / 8.2$
\end{codebox}

\begin{codebox}
    \Procname{$\proc{Bar}(x, y)$}
    \li $z \gets x + y$
    \li $w \gets x \times y$
    \li $q \gets (w^z) \bmod 870$
    \li \Return $9 \times q$
\end{codebox}

I det sidste eksempel udfører funktionen $\proc{Bar}$ en række regneoperationer, tildelinger og en modulo-operation. Hver af disse operationer regnes som en basisoperation, og funktionens tidskompleksitet er $\mathcal{O}(1)$, fordi den udfører et konstant antal basisoperationer. Her er en opdeling:

\begin{itemize}
    \item \textbf{Linje 1:} $z \gets x + y$
    \begin{itemize}
        \item \textbf{Regneoperation}: 1 (addition)
        \item \textbf{Tildeling}: 1 (tildeling til $z$)
    \end{itemize}
    \textbf{Operationer i alt}: 2
    
    \item \textbf{Linje 2:} $w \gets x \times y$
    \begin{itemize}
        \item \textbf{Regneoperation}: 1 (multiplikation)
        \item \textbf{Tildeling}: 1 (tildeling til $w$)
    \end{itemize}
    \textbf{Operationer i alt}: 2
    
    \item \textbf{Linje 3:} $q \gets (w^z) \bmod 870$
    \begin{itemize}
        \item \textbf{Regneoperationer}: 2 (én for potensopløftningen $w^z$ og én for modulo 870)
        \item \textbf{Tildeling}: 1 (tildeling til $q$)
    \end{itemize}
    \textbf{Operationer i alt}: 3
    
    \item \textbf{Linje 4:} \Return $9 \times q$
    \begin{itemize}
        \item \textbf{Regneoperation}: 1 (multiplikation med 9)
        \item \textbf{Tildeling (implicit)}: 1 (returnering af resultatet)
    \end{itemize}
    \textbf{Operationer i alt}: 2
\end{itemize}

Når vi lægger operationerne for hver linje sammen, får vi:
\[
2 + 2 + 3 + 2 = 9
\]

Det samlede antal basisoperationer er konstant (9 operationer), så dette er en $O(1)$-algoritme.

\subsection*{Sekventielle sætninger}

Når en algoritme indeholder flere sætninger, der udføres efter hinanden, kalder man dem »sekventielle sætninger« eller »på hinanden følgende sætninger«. Den samlede tidskompleksitet af sekventielle sætninger afhænger af de enkelte sætningers kompleksitet. Mere præcist er det leddet af højeste orden blandt de enkelte sætninger, der dominerer den samlede kompleksitet.

Se på en sekvens af sætninger:

\begin{align*}
    \text{sætning 1;} \\
    \text{sætning 2;} \\
    & \vdots \\
    \text{sætning k;}
\end{align*}

Den samlede tid for at udføre disse sætninger kan skrives som:

\[
\text{Samlet tid} = \text{tid(sætning 1)} + \text{tid(sætning 2)} + \cdots + \text{tid(sætning } k) 
\]

Hvis hver sætning er »simpel«, dvs. kun består af basisoperationer, tager hver af sætningerne konstant tid, og den samlede kompleksitet bliver $\mathcal{O}(1)$. Det gælder, når antallet af operationer i hver sætning ikke afhænger af inputstørrelsen -- det så vi i forrige afsnit.

I mere generelle tilfælde er det dog den højeste orden blandt de enkelte kompleksiteter, der bestemmer kompleksiteten af de sekventielle sætninger. Hvis fx én sætning har kompleksiteten $\mathcal{O}(n)$ og en anden har kompleksiteten $\mathcal{O}(n^2)$, er den samlede kompleksitet $\mathcal{O}(n^2)$, fordi leddet af højeste orden dominerer for tilstrækkeligt store inputstørrelser.

Som vi har set flere gange, bliver led af lavere orden ubetydelige i asymptotisk analyse, når inputstørrelsen vokser, og derfor medtages de ikke i den endelige angivelse af algoritmens kompleksitet. Når vi vurderer sekventielle sætninger, behøver vi altså kun at se på leddet med den højeste vækstrate for at vurdere den samlede tidskompleksitet.

\subsection*{Simple løkker}

Løkker er blandt de vigtigste faktorer for en algoritmes kompleksitet, fordi de ofte bestemmer, hvor mange gange en bestemt operation udføres. En simpel løkke, der kører fra 1 til $n$, har tidskompleksiteten $\mathcal{O}(n)$. Det skyldes, at løkken udføres $n$ gange, og hvis sekvensen af sætninger i løkken har kompleksiteten $\mathcal{O}(1)$, bliver løkkens samlede tidskompleksitet $n \times \mathcal{O}(1)$, hvilket kan reduceres til $\mathcal{O}(n)$. 

Den generelle form for en løkke kan skrives som:

\[
\begin{aligned}
& \text{for } i \text{ i } 1, 2, \ldots, N \text{ løkke} \\
& \quad \text{sekvens af sætninger} \\
& \text{slut løkke;}
\end{aligned}
\]

Her udføres løkken $N$ gange, hvilket betyder, at sekvensen af sætninger også udføres $N$ gange. Hvis vi antager, at hver af disse sætninger er en basisoperation med kompleksiteten $\mathcal{O}(1)$, er løkkens samlede tidskompleksitet $\mathcal{O}(N)$. 

Et illustrativt eksempel er en funktion, der udfører multiplikation uden at bruge den indbyggede multiplikationsoperator (\texttt{*}). I stedet bruger funktionen en løkke til gentagne gange at lægge den første operand til:

\begin{codebox}
    \Procname{$\proc{Mul2}(x, y)$}
    \li $result \gets 0$
    \li \For $i \gets 0$ \To $y - 1$ \Do
    \li     $result \gets result + x$
        \End
    \li \Return $result$
\end{codebox}

I denne funktion kører løkken $y$ gange og udfører en addition i konstant tid i hver iteration. Funktionens kompleksitet er derfor $\mathcal{O}(y)$. Tidskompleksiteten afhænger her direkte af størrelsen af inputtet $y$, hvilket viser, hvordan løkkegrænserne direkte påvirker en algoritmes ydelse.

Funktioner med simple for-løkker, der gennemløber hele inputtet, har generelt kompleksiteten $\mathcal{O}(n)$, hvor $n$ er inputstørrelsen. Det skyldes, at hvert element i inputtet behandles i rækkefølge, så der er en lineær sammenhæng mellem inputstørrelsen og antallet af udførte operationer.

Se på følgende:

\begin{codebox}
    \Procname{$\proc{Factorial}(n)$}
    \li $result \gets 1$
    \li \For $num \gets 1$ \To $n$ \Do
    \li     $result \gets result \times num$
        \End
    \li \Return $result$
\end{codebox}


Koden, der beregner fakultetet af et tal \( n \), består af en enkelt løkke, der gennemløber værdierne fra 1 til \( n \). I hver iteration udføres en multiplikation i konstant tid (result *= num). Løkken kører altså \( n \) gange, og hver iteration er en \(\mathcal{O}(1)\)-operation.

Den samlede tidskompleksitet for funktionen $\proc{factorial}$ er derfor $\mathcal{O}(n)$. Se nu på:

\begin{codebox}
    \Procname{$\proc{Factorial2}(n)$}
    \li $result \gets 1$
    \li $count \gets 0$
    \li \For $num \gets 1$ \To $n$ \Do
    \li     $result \gets result \times num$
    \li     $count \gets count + 1$
        \End
    \li \Return $result$
\end{codebox}

Funktionen $\proc{factorial2}$ minder meget om det forrige eksempel, hvor vi beregnede fakultetet af et tal \( n \). Den indeholder en enkelt løkke, der gennemløber værdierne fra 1 til \( n \). I hver iteration udføres to operationer i konstant tid: en multiplikation (result *= num) og en addition (count += 1). Da begge operationer er \(\mathcal{O}(1)\), og løkken kører \( n \) gange, er den samlede tidskompleksitet stadig \(\mathcal{O}(n)\).

Den ekstra tildeling (count += 1) ændrer ikke den \textit{asymptotiske} kompleksitet, da det stadig er en operation i konstant tid, der gentages \( n \) gange. Den samlede kompleksitet forbliver altså \(\mathcal{O}(n)\). Den ekstra operation kan dog påvirke funktionens ydelse i praksis, fordi den øger antallet af basisoperationer, der udføres i hver iteration.

\subsection*{If-then-else-sætninger}
For en \texttt{if-then-else}-struktur afhænger køretiden af omkostningen ved at evaluere betingelsen og af den gren, der udføres. Når vi analyserer det \emph{værste tilfælde}, bruger vi den langsomste gren plus omkostningen ved betingelsen. For et bestemt input medtages kun omkostningen ved den gren, der faktisk udføres.

\begin{itemize}
    \item \textbf{Hvis betingelsen er simpel}, fx en sammenligning af to tal, tager den konstant tid, $\mathcal{O}(1)$.
    \item \textbf{Hvis betingelsen er mere kompleks}, fx en test af, om et element findes i en stor datastruktur, afspejler tidskompleksiteten evalueringen af betingelsen.
\end{itemize}

Se på følgende pseudokode for en \texttt{if-then-else}-sætning:

\begin{codebox}
    \li \If cond \Then
    \li     \text{block 1 (sequence of statements)}
    \li \Else
    \li     \text{block 2 (sequence of statements)}
    \li \End
\end{codebox}


Her udføres enten blok 1 eller blok 2. Tiden i værste tilfælde er derfor den langsomste af de to muligheder:

\[
\max (\text{tid(blok 1)}, \text{tid(blok 2)})
\]

Hvis blok 1 tager $\mathcal{O}(1)$, og blok 2 tager $\mathcal{O}(N)$, er den samlede tidskompleksitet for \texttt{if-then-else}-sætningen $\mathcal{O}(N)$.

Kompleksiteten af betingelser afhænger af selve betingelsen. Tiden til at evaluere en betingelse kan være konstant, lineær eller endnu værre, afhængigt af hvilken test der udføres.

Se på følgende eksempler, der illustrerer dette nærmere:

\begin{codebox}
    \Procname{$\proc{Count\_ts}(a\_str)$}
    \li $count \gets 0$
    \li \For each $char$ in $a\_str$ \Do
    \li     \If $char = \text{'t'}$ \Then
    \li         $count \gets count + 1$
            \End
        \End
    \li \Return $count$
\end{codebox}


I dette eksempel brugte vi en \texttt{\If}-sætning til at teste, om ét tegn er lig med et andet, hvilket er en operation i konstant tid. Funktionen kører samlet set i lineær tid i forhold til længden af \texttt{a\_str}, fordi testen i betingelsen er simpel. Hvis vi lader $n = |\texttt{a\_str}|$, har funktionen tidskompleksiteten $\mathcal{O}(n)$.

Se nu på denne kode:

\begin{codebox}
    \Procname{$\proc{Count\_same\_ltrs}(a\_str, b\_str)$}
    \li $count \gets 0$
    \li \For each $char$ in $a\_str$ \Do
    \li     \If $char \in b\_str$ \Then
    \li         $count \gets count + 1$
            \End
        \End
    \li \Return $count$
\end{codebox}

Koden ligner $\proc{count\_ts}$, men testen i betingelsen, \texttt{char $\in$ b\_str}, er mere kompleks. Vi analyserer køretiden i værste tilfælde: En medlemskabstest kan blive nødt til at undersøge hvert tegn i \texttt{b\_str}. Det sker, når \texttt{char} ikke findes, fordi søgningen ikke kan returnere \texttt{False}, før alle tegn er kontrolleret. Store $\mathcal{O}$ betegner en øvre grænse; valget om at undersøge det værste tilfælde er en separat modelleringsbeslutning.

Lad $n = |\texttt{a\_str}|$ og $m = |\texttt{b\_str}|$. \texttt{\For}-løkken udfører $\mathcal{O}(n)$ iterationer, og hver iteration udfører en test i betingelsen på $\mathcal{O}(m)$. Da vi udfører en test på $\mathcal{O}(m)$ i alt $\mathcal{O}(n)$ gange, er funktionens samlede kompleksitet $\mathcal{O}(nm)$.

\subsection*{While-løkker}

While-løkker er grundlæggende konstruktioner til at skabe iterative processer i algoritmer. I modsætning til \texttt{\For}-løkker, der gennemløbes et forudbestemt antal gange, fortsætter en \texttt{\While}-løkke, så længe en given betingelse er sand. Denne fleksibilitet gør \texttt{\While}-løkker velegnede i situationer, hvor antallet af iterationer ikke kendes på forhånd, men afhænger af forhold under kørslen.

Strukturen af en \texttt{\While}-løkke kan skrives som:

\[
\text{så længe betingelsen er sand:}
\]
\[
\quad \text{sekvens af sætninger}
\]

\begin{itemize}
    \item Betingelsen kontrolleres før hver iteration. Hvis den er falsk fra starten, udføres løkkens krop slet ikke.
    \item Antallet af iterationer afhænger af, hvordan betingelsen udvikler sig og til sidst bliver falsk.
\end{itemize}

Tidskompleksiteten af en \texttt{\While}-løkke afhænger af kompleksiteten af testen i betingelsen og af det antal iterationer, der skal til, før betingelsen bliver falsk.

Se på følgende eksempel:

\begin{codebox}
    \Procname{$\proc{Factorial3}(n)$}
    \li $result \gets 1$
    \li \While $n > 0$ \Do
    \li     $result \gets result \times n$
    \li     $n \gets n - 1$
        \End
    \li \Return $result$
\end{codebox}

I dette eksempel fortsætter løkken, så længe \texttt{n > 0}. Da \texttt{n} formindskes med 1 i hver iteration, kører løkken $n$ gange, hvilket giver tidskompleksiteten $\mathcal{O}(n)$.

Se på et andet eksempel, der bruger en \texttt{\While}-løkke til at opdele en streng:

\begin{codebox}
    \Procname{$\proc{Char-Split}(a\_str)$}
    \li $result \gets [\,]$ \Comment{initialise an empty list}
    \li $index \gets 0$
    \li \While $len(a\_str) \neq len(result)$ \Do
    \li     $append(result, a\_str[index])$
    \li     $index \gets index + 1$
        \End
    \li \Return $result$
\end{codebox}

I denne funktion er $\texttt{len(a\_str)}$ og $\texttt{len(result)}$ operationer i konstant tid, og det samme gælder indeksering i strengen og tilføjelse til listen. While-løkken kører, indtil længden af $\texttt{result}$ svarer til længden af $\texttt{a\_str}$, hvilket giver $n$ iterationer, hvor $n$ er inputstrengens længde. Funktionens samlede tidskompleksitet er derfor $\mathcal{O}(n)$.

Se på dette eksempel:

\begin{codebox}
    \Procname{$\proc{Sum-Until-Threshold}(arr, threshold)$}
    \li $sum \gets 0$
    \li $i \gets 0$
    \li \While $sum < threshold$ \text{ and } $i < \attrib{arr}{length}$ \Do
    \li     $sum \gets sum + arr[i]$
    \li     $i \gets i + 1$
        \End
    \li \Return $sum$
\end{codebox}

Her kører løkken, indtil betingelsen \texttt{sum < threshold} ikke længere er sand, eller indtil alle elementer i arrayet \texttt{arr} er behandlet. Køretiden afhænger både af arrayets indhold og af værdien af \texttt{threshold}. I værste tilfælde undersøger løkken hvert element, så køretiden i værste tilfælde er $\Theta(n)$, hvor $n$ er arrayets længde.

Generelt er \texttt{while}-løkker stærke konstruktioner til at iterere over betingelser, der ikke kendes på forhånd. Nøglen til at analysere deres kompleksitet er at forstå, hvordan betingelsen udvikler sig i hver iteration, og at vurdere, hvor mange iterationer der skal til, før løkken stopper. En grundig analyse kræver, at man forstår både betingelsens udvikling og omkostningen ved hver iteration.

\subsection*{Indlejrede løkker}

Indlejrede løkker er en vigtig konstruktion i algoritmedesign, hvor én løkke kører inde i en anden. Denne type struktur har stor betydning for en algoritmes tidskompleksitet, fordi antallene af iterationer ganges med hinanden.



Se på den generelle form for en indlejret løkke:

\newpage

\[
\text{for } i \gets 1 \text{ til } n:
\]
\[
\quad \text{sekvens af sætninger}
\]
\[
\quad \text{for } j \gets 1 \text{ til } m:
\]
\[
\quad\quad \text{sekvens af sætninger}
\]

Ved indlejrede løkker gennemløbes den indre løkke fuldt ud for hver iteration af den ydre løkke. Det betyder, at det samlede antal iterationer er produktet af de enkelte løkkers antal iterationer. Hvis den ydre løkke kører $n$ gange og den indre løkke $m$ gange, udføres den indre løkkes krop i alt $n \times m$ gange. Hvis både $m$ og $n$ er proportionale med inputstørrelsen, er den indlejrede løkkes tidskompleksitet $\mathcal{O}(n \times m)$. 

Hvis både den indre og den ydre løkke gennemløber det samme interval, fx fra 1 til $n$, er den samlede tidskompleksitet for strukturen af indlejrede løkker $\mathcal{O}(n^2)$. Det er ofte tilfældet, når man udfører operationer på en todimensional struktur som en matrix, hvor hver række og søjle skal gennemløbes.

For eksempel:

\begin{codebox}
    \Procname{$\proc{Sum-Matrix}(matrix)$}
    \li $sum \gets 0$
    \li \For $i \gets 0$ \To $\attrib{matrix}{rows} - 1$ \Do
    \li     \For $j \gets 0$ \To $\attrib{matrix}{cols} - 1$ \Do
    \li         $sum \gets sum + matrix[i][j]$
            \End
        \End
    \li \Return $sum$
\end{codebox}

I denne pseudokode gennemløber den ydre løkke hver række i matricen, mens den indre løkke gennemløber hver søjle. Hvis matricen har $n$ rækker og $m$ søjler, er tidskompleksiteten af den indlejrede løkke $\mathcal{O}(n \times m)$. For en kvadratisk matrix, hvor $n = m$, bliver kompleksiteten $\mathcal{O}(n^2)$.

I nogle tilfælde afhænger antallet af iterationer i den indre løkke af den ydre løkkes indeks. Et almindeligt eksempel er en trekantet løkkestruktur, hvor den indre løkke kører færre gange, efterhånden som den ydre løkke skrider frem.

Se på følgende pseudokode:

\begin{codebox}
    \Procname{$\proc{Triangular-Sum}(n)$}
    \li $sum \gets 0$
    \li \For $i \gets 1$ \To $n$ \Do
    \li     \For $j \gets i$ \To $n$ \Do
    \li         $sum \gets sum + (i + j)$
            \End
        \End
    \li \Return $sum$
\end{codebox}

Her starter den indre løkke ved den aktuelle værdi af $i$ og kører til $n$. Det betyder, at den indre løkke i første iteration kører $n$ gange, i anden iteration $n-1$ gange og så videre, indtil den kun kører én gang. Det samlede antal iterationer for de indlejrede løkker kan beregnes som (se Bilag \ref{app:appendix01}):

\[
\sum_{i=1}^{n} (n - i + 1) = \frac{n (n + 1)}{2}
\]

Tidskompleksiteten af denne indlejrede løkke er altså $\mathcal{O}(n^2)$. Den trekantede struktur reducerer det samlede antal iterationer sammenlignet med en fuld $\mathcal{O}(n^2)$-løkke, men kompleksiteten er stadig kvadratisk.

Det samme princip gælder for mere end to indlejrede løkker. Hvis der er $k$ indlejrede løkker, der hver kører $n$ gange, er den samlede tidskompleksitet $\mathcal{O}(n^k)$.

Når man analyserer indlejrede løkker, er det afgørende at forstå hver løkkes interval og grænser. Hvis grænserne er dynamiske og afhænger af værdien af andre variabler, kan kompleksiteten være sværere at bestemme, og det kan kræve en grundigere analyse af, hvordan grænserne ændrer sig under udførelsen.

Indlejrede løkker bruges ofte til problemer med flerdimensionale datastrukturer eller til at sammenligne hvert element i en samling med alle andre elementer (fx i sorteringsalgoritmer som $\proc{Bubble Sort}$, der har tidskompleksiteten $\mathcal{O}(n^2)$ på grund af sin struktur med indlejrede løkker).

\subsection*{Rekursion}
Denne bog og de kurser, den understøtter, behandler ikke rekursion i dybden. Vi omtaler kort rekursion her i forbindelse med konstruktionsanalyse for fuldstændighedens skyld. Vi har allerede set rekursion i forbindelse med binær søgning. Rekursion er et stærkt værktøj i algoritmedesign, hvor funktioner kalder sig selv og løser problemer ved at opdele dem i mindre delproblemer. Rekursive algoritmer er særligt effektive til problemer med en naturlig hierarkisk eller gentagen struktur, fx gennemløb af træer, del og hersk-algoritmer eller beregning af fakulteter.

En central del af enhver rekursiv funktion er \textbf{basistilfældet} -- den betingelse, hvor funktionen holder op med at kalde sig selv, så man undgår uendelig rekursion.

Se på den generelle struktur af en rekursiv funktion:

\begin{itemize}
    \item \textbf{Basistilfælde}: Den stopbetingelse, der afslutter rekursionen.
    \item \textbf{Rekursivt tilfælde}: Den del af funktionen, der kalder sig selv med et ændret input og dermed nærmer sig basistilfældet.
\end{itemize}

Et klassisk eksempel på en rekursiv funktion er beregningen af et fakultet:

\[
\texttt{factorial}(n) = 
\begin{cases}
1 & \text{hvis } n = 0 \\
n \times \texttt{factorial}(n-1) & \text{hvis } n > 0
\end{cases}
\]

I pseudokode kan fakultetsfunktionen skrives som:

\begin{codebox}
    \Procname{$\proc{Factorial4}(n)$}
    \li \If $n == 0$ \Then
    \li     \Return $1$
        \End
    \li \Return $n \times \proc{Factorial}(n-1)$
\end{codebox}

I denne pseudokode:
\begin{itemize}
    \item \textbf{Basistilfældet} indtræffer, når $n == 0$, og returnerer $1$.
    \item \textbf{Det rekursive tilfælde} ganger $n$ med resultatet af $\proc{Factorial4}(n-1)$.
\end{itemize}

Tidskompleksiteten af en rekursiv funktion afhænger af, hvor mange gange funktionen kalder sig selv, og hvor meget arbejde der udføres på hvert rekursionsniveau. For fakultetsfunktionen gælder:

\begin{itemize}
    \item Funktionen foretager ét rekursivt kald i hvert trin og formindsker $n$ med $1$ hver gang.
    \item Rekursionsdybden er $n$, og hvert kald udfører $\mathcal{O}(1)$ arbejde (en multiplikation).
\end{itemize}

Fakultetsfunktionens tidskompleksitet er altså $\mathcal{O}(n)$.

Til at bestemme tidskompleksiteten af en rekursiv algoritme kan vi bruge rekursionsligninger. En rekursionsligning udtrykker en rekursiv funktions tidskompleksitet ved hjælp af dens delproblemer. Rekursionsligningen for fakultetsfunktionen kan fx skrives som:

\[
T(n) = T(n-1) + \mathcal{O}(1)
\]

Det betyder, at tiden til at beregne $\texttt{factorial}(n)$ er lig med tiden til at beregne $\texttt{factorial}(n-1)$ plus en konstant mængde arbejde, hvilket giver kompleksiteten $\mathcal{O}(n)$.

Et andet velkendt eksempel på en rekursiv funktion er Fibonacci-følgen:

\[
\texttt{fibonacci}(n) = 
\begin{cases}
0 & \text{hvis } n = 0 \\
1 & \text{hvis } n = 1 \\
\texttt{fibonacci}(n-1) + \texttt{fibonacci}(n-2) & \text{hvis } n > 1
\end{cases}
\]

Den tilsvarende pseudokode er:

\begin{codebox}
    \Procname{$\proc{Naive-Fibonacci}(n)$}
    \li \If $n == 0$ \Then
    \li     \Return $0$
        \End
    \li \If $n == 1$ \Then
    \li     \Return $1$
        \End
    \li \Return $\proc{Fibonacci}(n-1) + \proc{Fibonacci}(n-2)$
\end{codebox}

I dette eksempel:
\begin{itemize}
    \item Funktionen kalder sig selv to gange i det rekursive tilfælde, hvilket skaber en forgrenet struktur.
    \item \textbf{Tidskompleksiteten} er $\mathcal{O}(2^n)$, som vokser eksponentielt, fordi de samme delproblemer løses mange gange. Det viser, hvor ineffektiv den naive rekursive implementering af Fibonacci er.
\end{itemize}

For at gøre rekursive funktioner, der løser overlappende delproblemer som i Fibonacci-følgen, mere effektive kan man bruge \textbf{memoisering}. Memoisering går ud på at gemme resultaterne af delproblemer for at undgå overflødige beregninger.

Hvis man fx anvender memoisering på Fibonacci-funktionen, ændres tidskompleksiteten fra $\mathcal{O}(2^n)$ til $\mathcal{O}(n)$, fordi hvert delproblem kun løses én gang.

Rekursion kan ofte erstattes af iteration, især for problemer, hvor rekursionsdybden kan blive et problem. Fakultetsfunktionen kan fx omskrives iterativt, som vi så tidligere. Her er den iterative version af fakultetsfunktionen:

\begin{codebox}
    \Procname{$\proc{Factorial-Iterative}(n)$}
    \li $result \gets 1$
    \li \For $i \gets 1$ \To $n$ \Do
    \li     $result \gets result \times i$
        \End
    \li \Return $result$
\end{codebox}

Den iterative fremgangsmåde har samme tidskompleksitet, $\mathcal{O}(n)$, som den rekursive version, men undgår omkostningerne ved rekursive kald og er ikke begrænset af systemets stakstørrelse.

\textbf{Halerekursion} er en særlig form for rekursion, hvor det rekursive kald er den sidste operation i funktionen. I nogle sprog kan compileren optimere halerekursive funktioner til iteration, så de ikke bruger ekstra plads på stakken. Se på følgende halerekursive fakultetsfunktion:

\begin{codebox}
    \Procname{$\proc{Factorial-Tail}(n, accumulator)$}
    \li \If $n == 0$ \Then
    \li     \Return $accumulator$
        \End
    \li \Return $\proc{Factorial-Tail}(n-1, accumulator \times n)$
\end{codebox}

Her står det rekursive kald i \textbf{haleposition}, hvilket betyder, at der ikke skal foretages yderligere beregninger, efter at kaldet har returneret. Halerekursion mindsker risikoen for \textit{stakoverløb} (på engelsk \emph{stack overflow}) og bliver ofte optimeret til en løkke af compileren.

Rekursion er et alsidigt værktøj i algoritmedesign, der giver elegante løsninger på komplekse problemer ved at opdele dem i mindre og mere håndterbare delproblemer. Den indebærer dog afvejninger med hensyn til plads- og tidskompleksitet, især ved dyb eller eksponentiel rekursion.

Analyse af rekursive algoritmer kræver, at man forstår de rekursionsligninger, der beskriver deres opførsel, identificerer basistilfældene og vurderer, hvor mange gange funktionen bliver kaldt. Rekursion kan give klar og kortfattet kode, men ofte er en iterativ fremgangsmåde eller optimeringsteknikker som memoisering og optimering af halekald nødvendige for at forbedre effektiviteten.

\section{Analyse af løkker}\label{sec:loop-analysis}
I forrige afsnit har vi allerede diskuteret løkkers tidskompleksitet. I dette afsnit analyserer vi løkker mere detaljeret, herunder indlejrede løkker, løkker med variable grænser og løkker med betingede udgange.

Når der i denne bog bedes om en klassifikation med store $\mathcal{O}$, menes der normalt den mest informative standardgrænse opadtil blandt de almindelige vækstklasser, der er i spil. Selv om $\mathcal{O}(n^2)$ er en gyldig øvre grænse for $f(n)=4n+3$, er $\mathcal{O}(n)$ altså mere informativ. Strengt taget udtrykker store $\mathcal{O}$ alene ikke nødvendigvis, at grænsen er stram; $f(n)=\Theta(n)$ angiver både en øvre og en nedre asymptotisk grænse, der passer sammen.

\subsection*{Teknikker til praktisk analyse af løkker}

Nu hvor vi har set de forskellige typer løkker og deres generelle bidrag til kompleksiteten, er det tid til at gå i dybden med praktiske teknikker til at analysere disse konstruktioner i virkelige situationer. Dette afsnit giver konkrete metoder, som læseren kan bruge til systematisk at vurdere løkkerne i sine algoritmer, med fokus på både de teoretiske og de praktiske sider af analysen.

\subsubsection*{Bestemmelse af løkkegrænser}
For at bestemme en løkkes tidskompleksitet begynder man med at fastlægge grænserne for gennemløbene. Det indebærer, at man identificerer startpunktet, slutpunktet og mønstret for, hvordan løkkevariablerne ændres. Løkkegrænserne fastlægger den øvre grænse for antallet af iterationer, så vi kan udlede det værste tilfælde for, hvor mange gange en løkke udføres.

Se fx på en løkke af formen:
\[ 
\text{for } i \gets a \text{ til } b \text{ med skridt } c:
\]
Antallet af iterationer i denne løkke kan beregnes som:
\[
\left(\frac{b - a}{c}\right) + 1
\]
hvor $a$ er startværdien, $b$ er slutværdien, og $c$ er skridtlængden. En analyse af disse grænser hjælper med at bestemme løkkens samlede kompleksitet.

De følgende eksempler illustrerer begrebet.

\begin{example}

    \begin{codebox}
        \Procname{$\proc{Loop-Constant-Increment}(n)$}
        \li \For $i \gets 1$ \To $n$ \Do
        \li     \Comment{Perform some $\mathcal{O}(1)$ operation}
            \End
    \end{codebox}

I dette eksempel kører løkken fra $a$ til $b$ med skridtlængden $1$. Antallet af iterationer er:

\[
\left( \frac{b - a}{1} \right) + 1 = (b - a + 1) = (n - 1 + 1) = n 
\]

I dette eksempel starter $i$ ved $1$ og øges med $1$ indtil $n$. Løkken udføres $n$ gange, hvilket giver tidskompleksiteten $\mathcal{O}(n)$.

\end{example}

Eksemplet ovenfor viser direkte, hvordan formlen for antallet af iterationer bruges, når skridtlængden er konstant. Her er $a = 1$, $b = n$ og $c = 1$, hvilket giver $n$ iterationer.

\begin{example}

    \begin{codebox}
        \Procname{$\proc{Loop-Geometric-Growth}(n)$}
        \li $i \gets 1$
        \li \While $i \leq n$ \Do
        \li     \Comment{Perform some $\mathcal{O}(1)$ operation}
        \li     $i \gets 2 \times i$
            \End
    \end{codebox}
    
    I dette eksempel starter $i$ ved $1$ og fordobles i hver iteration, indtil den overstiger $n$. Antallet af iterationer kan bestemmes således:
    
    Værdien af $i$ udvikler sig som:
    
    \[
    1, 2, 4, 8, \dots, 2^k
    \]
    
    hvor $2^k \leq n$. For at bestemme antallet af iterationer løser vi:
    
    \[
    2^k \leq n \implies k \leq \log_2(n)
    \]
    
    Antallet af iterationer er altså cirka $\log_2(n)$, hvilket giver tidskompleksiteten:
    
    \[
    \mathcal{O}(\log n)
    \]
    
    I dette eksempel er ændringsmønstret multiplikativt i stedet for additivt. Selv om den tidligere formel kun gælder direkte for additive skridt, gælder princippet om at analysere løkkegrænserne for at bestemme antallet af iterationer stadig. Her bestemmer vi, hvor mange gange vi kan fordoble $i$, før den overstiger $n$, hvilket giver en logaritmisk kompleksitet.

\end{example}

\begin{example}

    \begin{codebox}
        \Procname{$\proc{Loop-Decrement}(n)$}
        \li $i \gets n$
        \li \While $i > 0$ \Do
        \li     \Comment{Perform some $\mathcal{O}(1)$ operation}
        \li     $i \gets i - 3$
            \End
    \end{codebox}

    I dette eksempel starter $i$ ved $n$ og formindskes med $3$ hver gang, indtil den bliver mindre end eller lig med $0$. Antallet af iterationer er givet ved:

    \[
    i: n, (n - 3), (n - 6), \dots
    \]

    Løkken fortsætter, indtil $i \leq 0$. For at bestemme antallet af iterationer skal vi se på, hvor mange gange vi kan trække $3$ fra $n$:

    \[
    \frac{n}{3}
    \]

    Antallet af iterationer er altså cirka:

    \[
    \left( \frac{n}{3} \right)
    \]

    Løkkens samlede tidskompleksitet er:

    \[
    \mathcal{O}(n)
    \] 

    I dette eksempel er skridtlængden $-3$. Med den generelle formel fra tidligere har vi $a = n$, $b = 0$ og $c = -3$. Det giver et antal iterationer på cirka $n / 3$. Princippet om at fastlægge grænserne og bruge formlen gælder altså direkte her og viser, hvordan formindskelsen påvirker antallet af iterationer.

\end{example}


\subsubsection*{Kompleksitet gennem summation}
I algoritmer med flere løkker efter hinanden er den samlede kompleksitet summen af de enkelte løkkers kompleksitet. Når løkker er indlejrede, er kompleksiteten derimod produktet af de enkelte kompleksiteter. To løkker efter hinanden, der hver har kompleksiteten $\mathcal{O}(n)$, giver fx stadig $\mathcal{O}(n)$, mens en indlejret struktur giver $\mathcal{O}(n^2)$. Det er afgørende at forstå denne forskel.

Ved indlejrede løkker er det ofte nødvendigt at bestemme det samlede antal operationer ved hjælp af summation. Hvis den indre løkke afhænger af den ydre løkkes værdi, er det nyttigt at opskrive den sum, der angiver antallet af iterationer for alle løkkerne. For eksempel:
\[
\text{for } i \gets 1 \text{ til } n:
\]
\[
\quad \text{for } j \gets 1 \text{ til } i:
\]
Det samlede antal iterationer er givet ved summen (se Bilag \ref{app:appendix01}):
\[
\sum_{i=1}^{n} i = \frac{n(n + 1)}{2} = \mathcal{O}(n^2)
\]
Denne type analyse er særligt nyttig for indlejrede løkker med afhængige intervaller. De følgende eksempler illustrerer begrebet.

\begin{example}

    \begin{codebox}
        \Procname{$\proc{Linear-Sum-Large-Constant}(n)$}
        \li \For $i \gets 1$ \To $n$ \Do
        \li     \For $j \gets 1$ \To $10^9$ \Do
        \li         \Comment{Perform some $\mathcal{O}(1)$ operation}
                \End
            \End
    \end{codebox}

    I dette eksempel kører den ydre løkke fra $1$ til $n$, mens den indre løkke kører et konstant antal gange (fra $1$ til $10^9$). Det samlede antal iterationer er:

    \[
    \sum_{i=1}^{n} 10^9 = 10^9 \times n = \mathcal{O}(n)
    \]

    Selv om den indre løkke har et meget stort konstant antal iterationer, er den samlede kompleksitet stadig lineær i $n$, fordi den konstante faktor ikke påvirker den asymptotiske vækst.
\end{example}

\begin{example}

    \begin{codebox}
        \Procname{$\proc{Quadratic-Sum}(n)$}
        \li \For $i \gets 1$ \To $n$ \Do
        \li     \For $j \gets 1$ \To $i$ \Do
        \li         \Comment{Perform some $\mathcal{O}(1)$ operation}
                \End
            \End
    \end{codebox}

I dette eksempel kører den ydre løkke fra $1$ til $n$, og den indre løkke kører fra $1$ til $i$. Det samlede antal iterationer er:

    \[
        \sum_{i=1}^n \sum_{j=1}^i 1 = \sum_{i=1}^{n} i = \frac{n(n + 1)}{2} = \mathcal{O}(n^2)
    \]

I dette udtryk gælder:
\begin{itemize}
    \item Den ydre sum, $\sum_{i=1}^n$, svarer til den ydre løkkes iterationer over $i$.
    \item Den indre sum, $\sum_{j=1}^i$, svarer til den indre løkkes iterationer for hvert fast $i$.
\end{itemize}

Tallet 1 i den indre sum svarer til, at der udføres en operation i konstant tid i hver iteration af den indre løkke. Udtrykket tæller hver iteration af den indre løkke enkeltvis for alle par af $i$ og $j$, der opfylder $1 \leq j \leq i$. Dobbeltsummen svarer til $\sum_{i=1}^n i$, fordi den indre løkke for hvert $i$ kører præcis $i$ gange. For at se, hvorfor den enkelte sum svarer til dobbeltsummen, kan vi prøve at implementere den enkelte sum:

\begin{codebox}
    \Procname{$\proc{Sum-Up-To}(n)$}
    \li $total \gets 0$
    \li \For $i \gets 1$ \To $n$ \Do
    \li     $total \gets total + i$
        \End
    \li \Return $total$
\end{codebox}

I denne kode:

\begin{itemize}
    \item Den ydre løkke kører fra 1 til $n$.
    \item I hver iteration lægges den aktuelle værdi af $i$ til \texttt{total}.
    \item Når løkken er færdig, indeholder \texttt{total} summen $1+2+\cdots+n$, som netop er $\sum_{i=1}^n i$.
\end{itemize}

Summen \( \sum_{i=1}^n i = \frac{n(n + 1)}{2} \) vokser proportionalt med \( n^2 \), men tidskompleksiteten af funktionen $\proc{Sum-Up-To(n)}$ er stadig lineær, \( \mathcal{O}(n) \), fordi den udfører en operation i konstant tid i hver af sine \( n \) iterationer. Det viser, at en algoritmes output kan have en talværdi, der vokser kvadratisk med \( n \), mens dens tidskompleksitet afhænger af antallet af operationer, som vokser lineært.

\end{example}

\begin{example}

    \begin{codebox}
        \Procname{$\proc{Logarithmic-Growth}(n)$}
        \li \For $i \gets 1$ \To $n$ \Do
        \li     $j \gets 1$
        \li     \While $j < n$ \Do
        \li         \Comment{Perform some $\mathcal{O}(1)$ operation}
        \li         $j \gets 2 \times j$
                \End
            \End
    \end{codebox}

    I dette eksempel kører den ydre løkke fra $1$ til $n$, mens den indre løkke starter ved $1$ og fordobler $j$, indtil den når eller overstiger $n$. Det samlede antal iterationer af den indre løkke over alle iterationer af den ydre løkke er:

    \[
    \sum_{i=1}^{n} \lceil\log_2 n\rceil = n\lceil\log_2 n\rceil = \mathcal{O}(n \log n)
    \]

    Den samlede tidskompleksitet for denne indlejrede struktur er altså $\mathcal{O}(n \log n)$.
\end{example}



\subsubsection*{Kompleksitet i værste, bedste og gennemsnitligt tilfælde}
Ikke alle løkker kører lige mange gange i alle situationer. Opførslen i bedste, værste og gennemsnitligt tilfælde kan variere afhængigt af den betingelse, der styrer, hvornår løkken stopper. Hvis en løkke fx afsluttes tidligt, fordi en betingelse er opfyldt, kan kompleksiteten i bedste tilfælde være betydeligt mindre end i værste tilfælde. Det er vigtigt at forstå disse tilfælde, og hvordan de påvirker køretiden, i virkelige applikationer, hvor inputtet kan variere.

\begin{example}

    \begin{codebox}
        \Procname{$\proc{Find-Target}(arr, target)$}
        \li $n \gets \attrib{arr}{length}$
        \li \For $i \gets 1$ \To $n$ \Do
        \li     \If $arr[i] == target$ \Then
        \li         \Return $i$ \Comment{Return the index of the target if found}
                \End
            \End
    \end{codebox}

    I dette eksempel er målet at finde et bestemt element (\emph{target}) i et array af længde $n$. Hvor mange iterationer løkken udfører, afhænger af elementets position:

    \begin{itemize}
        \item \textbf{Kompleksitet i bedste tilfælde:} Elementet findes på den første position i arrayet. Løkken kører kun én gang, hvilket giver kompleksiteten $\boldsymbol{\Theta(1)}$.\footnote{Man kunne tro, at den nedre grænse er $\Omega(1)$. I bedste tilfælde, hvor elementet står på den første position, udfører algoritmen en konstant mængde arbejde. Tidskompleksiteten er både $\boldsymbol{\mathcal{O}(1)}$ (øvre grænse) og $\boldsymbol{\Omega(1)}$ (nedre grænse), hvilket betyder, at den er $\boldsymbol{\Theta(1)}$ (en stram grænse). Det er mere præcist at bruge $\boldsymbol{\Theta(1)}$, fordi det angiver, at tiden er begrænset af en konstant både opadtil og nedadtil.}
        \item \textbf{Kompleksitet i værste tilfælde:} Elementet findes ikke i arrayet eller står på den sidste position. Løkken kører $n$ gange, hvilket giver kompleksiteten $\boldsymbol{\mathcal{O}(n)}$.
        \item \textbf{Kompleksitet i gennemsnitligt tilfælde:} I gennemsnit kan elementet stå et sted midt i arrayet. Hvis vi antager en ligefordeling, er det gennemsnitlige antal iterationer cirka $n/2$, hvilket giver kompleksiteten $\boldsymbol{\Theta(n)}$ i gennemsnitligt tilfælde.
    \end{itemize}    

    Eksemplet viser, hvordan den samme løkke kan opføre sig forskelligt afhængigt af inputtet, så der er forskellige kompleksiteter i bedste, værste og gennemsnitligt tilfælde.
\end{example}

\begin{example}


Se på den iterative binære søgealgoritme, som vi omtalte tidligere i kapitlet, anvendt på et sorteret array af længde $n$. Binær søgning er en del og hersk-algoritme, der virker ved gentagne gange at halvere søgeintervallet.

\begin{codebox}
    \Procname{$\proc{Binary-Search-Iterative}(arr, key)$}
    \li $low \gets 0$ \CommentAt{7}{Initialise low index}
    \li $high \gets arr.length -1$ \CommentAt{7}{Initialise high index}
    \li \While $low \leq high$ \Do
    \li     $mid \gets \floor{(low + high) / 2}$ \CommentAt{7}{Calculate mid index}
    \li     \If $arr[mid] == key$ \Then
    \li         \Return $mid$ \CommentAt{7}{Key found at mid index}
    \li     \ElseIf $arr[mid] < key$ \Then
    \li         $low \gets mid + 1$ \CommentAt{7}{Update low index}
    \li     \Else
    \li         $high \gets mid - 1$ \CommentAt{7}{Update high index}
        \End
        \End
    \li \Return $-1$ \CommentAt{7}{Key not found}
\end{codebox}

\begin{itemize}
    \item \textbf{Kompleksitet i bedste tilfælde: } 
    I bedste tilfælde står det søgte element midt i arrayet og findes ved den allerførste sammenligning. Da det kun kræver én operation, er kompleksiteten $\boldsymbol{\Theta(1)}$. Vi bruger denne notation, fordi den giver en stram grænse -- algoritmen bruger altid konstant tid i bedste tilfælde.

    \item \textbf{Kompleksitet i værste tilfælde: }
    I værste tilfælde bliver algoritmen ved med at halvere søgeintervallet, indtil der kun er én kandidat tilbage, eller intervallet er tomt. Antallet af iterationer vokser logaritmisk, så køretiden i værste tilfælde er $\boldsymbol{\Theta(\log n)}$.

    \item \textbf{Kompleksitet i gennemsnitligt tilfælde: }  
    Betinget af, at søgningen lykkes, og at det søgte element med lige stor sandsynlighed står på hver position i arrayet, er det forventede antal sammenligninger proportionalt med $\log n$. Under denne model er den forventede køretid derfor $\boldsymbol{\Theta(\log n)}$. En model, der også tillader mislykkede søgninger, må angive deres sandsynlighed, men giver samme asymptotiske orden under almindelige antagelser om en fast sandsynlighed.
\end{itemize}
\end{example}

\begin{example}

    Se på den naive strengsøgningsalgoritme (\emph{Naïve String Matching}) anvendt på en tekststreng $T$ af længde $n$ og en mønsterstreng $P$ af længde $m$. Målet er at finde alle forekomster af $P$ i $T$ ved at teste for et match ved hver mulig forskydning.
    
    \begin{codebox}
        \Procname{$\proc{Naive-String-Matching}(T, P)$}
        \li $n \gets \attrib{T}{length}$
        \li $m \gets \attrib{P}{length}$
        \li \For $s \gets 0$ \To $n - m$ \Do
        \li     $i \gets 0$
        \li     \While $i < m$ and $T[s + i] == P[i]$ \Do
        \li         $i \gets i + 1$
                \End
        \li     \If $i == m$ \Then
        \li         \textbf{report} $s$ \Comment{Pattern found at position $s$}
                \End
            \End
    \end{codebox}
    
    \begin{itemize}
        \item \textbf{Kompleksitet i bedste tilfælde: }  
        I bedste tilfælde fejler den første tegnsammenligning ved hver forskydning. Algoritmen tester stadig alle $n-m+1$ forskydninger, så køretiden i bedste tilfælde er $\boldsymbol{\Theta(n-m+1)}$ (og dermed $\Theta(n)$, når $m\leq cn$ for et fast $c<1$).
    
        \item \textbf{Kompleksitet i værste tilfælde: }
        I værste tilfælde indeholder teksten og mønstret gentagne sekvenser, der tvinger algoritmen til at foretage $m$ sammenligninger ved stort set hver forskydning. Køretiden i værste tilfælde er derfor $\boldsymbol{\Theta((n-m+1)m)}$, som normalt skrives $\Theta(nm)$, når $m$ ikke ligger så tæt på $n$, at det ændrer den første faktor.
    
        \item \textbf{Kompleksitet i gennemsnitligt tilfælde: } 
        Under en model med uafhængige tegn over et fast alfabet med mindst to symboler er det forventede antal sammenligninger pr. forskydning begrænset af en konstant. Den forventede køretid er derfor $\boldsymbol{\Theta(n-m+1)}$, som ofte forenkles til $\Theta(n)$ under samme konvention som ovenfor.
    \end{itemize}
   
    
    Eksemplet viser også, at bedste, værste og gennemsnitligt tilfælde beskriver forskellige køretidsfunktioner, mens $\mathcal{O}$, $\Omega$ og $\Theta$ beskriver grænser for den funktion, der analyseres.
    
    \end{example}

\subsection*{Eksempler på analyse af løkker}
Vi giver nu en række eksempler, der illustrerer, hvordan man analyserer løkker.

\begin{example}

    Se på følgende algoritme:
    
    \begin{codebox}
        \Procname{$\proc{Loop1}(n)$}
        \li $i \gets 1$
        \li \While $i \leq n$ \Do
        \li     $i \gets 3 \times i$
            \End
    \end{codebox}
    
    I hver iteration af \texttt{\While}-løkken ganges variablen $i$ med $3$. Vi vil bestemme, hvor mange gange løkken udføres, før $i$ overstiger $n$.
    
    Lad $k$ være det antal iterationer, løkken udfører.
    
    Til at begynde med er $i_0 = 1$.
    
    Efter første iteration:
    \[ i_1 = 3 \times i_0 = 3 \times 1 = 3 \]
    
    Efter anden iteration:
    \[ i_2 = 3 \times i_1 = 3 \times 3 = 3^2 = 9 \]
    
    Efter $k$ iterationer:
    \[ i_k = 3^k \times i_0 = 3^k \]
    
    Løkken fortsætter, så længe $i_k \leq n$. Betingelsen for, at løkken stopper, er derfor:
    \[ 3^k > n \]
    
    Vi tager logaritmen med grundtal $3$ på begge sider:
    \[ k > \log_3 n \]
    
    Da $k$ skal være et helt tal, er antallet af iterationer:
    \[ k = \lfloor \log_3 n \rfloor + 1 \]
    
    \textbf{Konklusion:} Løkken kører $\lfloor \log_3 n \rfloor + 1$ gange, så algoritmens tidskompleksitet er $\boldsymbol{\mathcal{O}(\log n)}$.
    
\end{example}
    
 \begin{example}

    Se på følgende algoritme:
    
    \begin{codebox}
        \Procname{$\proc{Loop2}(n)$}
        \li $s \gets 1$
        \li \For $i \gets 1$ \To $n$ \Do
        \li     \For $j \gets 1$ \To $n$ \Do
        \li         $s \gets s + 1$
                \End
            \End
    \end{codebox}
    
    
    Vi har to indlejrede \texttt{\For}-løkker:
    
    \begin{itemize}
        \item Den ydre løkke kører fra $i = 1$ til $i = n$ og udfører $n$ iterationer.
        \item Den indre løkke kører fra $j = 1$ til $j = n$ og udfører ligeledes $n$ iterationer for hver iteration af den ydre løkke.
    \end{itemize}
    
    I hver iteration af den indre løkke udføres en operation i konstant tid: $s \gets s + 1$.
    
    \textbf{Samlet antal iterationer:}
    
    \begin{itemize}
        \item Den indre løkke kører $n$ gange for hver af den ydre løkkes $n$ iterationer.
        \item Det samlede antal gange, $s \gets s + 1$ udføres, er derfor:
        \[ n \times n = n^2 \]
    \end{itemize}
    
    \textbf{Konklusion:} Algoritmen udfører omkring $n^2$ operationer, så tidskompleksiteten er $\boldsymbol{\mathcal{O}(n^2)}$.
        
\end{example}

\begin{example}

    Se på følgende algoritme:
    
    \begin{codebox}
        \Procname{$\proc{Loop3}(n)$}
        \li $i \gets 1$
        \li \While $i \leq n$ \Do
        \li     $i \gets 2 \times i$
            \End
    \end{codebox}
    
    I hver iteration af \texttt{\While}-løkken ganges variablen $i$ med $2$. Vi skal bestemme, hvor mange gange løkken udføres, før $i$ overstiger $n$.
    
    Lad $k$ være det antal iterationer, løkken udfører.
    
    \begin{itemize}
        \item \textbf{Initialisering:} $i_0 = 1$.
        \item \textbf{Opdatering i hver iteration:} Efter hver iteration opdateres $i$ som:
        \[ i_k = 2^k \times i_0 = 2^k \]
        \item \textbf{Stopbetingelse:} Løkken fortsætter, så længe $i_k \leq n$. Altså:
        \[ 2^k \leq n \]
        \item \textbf{Løsning med hensyn til $k$:} Vi tager logaritmen med grundtal $2$ på begge sider:
        \[ k \leq \log_2 n \]
        \item \textbf{Antal iterationer:} Da $k$ skal være et helt tal, udføres løkken:
        \[ k = \lfloor \log_2 n \rfloor + 1 \]
    \end{itemize}
    
    \textbf{Konklusion:} Løkken kører $\lfloor \log_2 n \rfloor + 1$ gange, så algoritmens tidskompleksitet er $\boldsymbol{\mathcal{O}(\log n)}$.
            
\end{example}
            
\begin{example}

    Se på følgende algoritme:
    
    \begin{codebox}
        \Procname{$\proc{Loop4}(n)$}
        \li $i \gets 1$
        \li \While $i \leq n \times n$ \Do
        \li     $i \gets 3 \times i$
            \End
    \end{codebox}
    
    I hver iteration af \texttt{\While}-løkken ganges variablen $i$ med $3$. Vi skal bestemme, hvor mange gange løkken udføres, før $i$ overstiger $n^2$.
    
    Lad $k$ være det antal iterationer, løkken udfører.
    
    \begin{itemize}
        \item \textbf{Initialisering:} $i_0 = 1$.
        \item \textbf{Opdatering i hver iteration:} Efter hver iteration:
        \[ i_k = 3^k \times i_0 = 3^k \]
        \item \textbf{Stopbetingelse:} Løkken fortsætter, så længe $i_k \leq n^2$. Altså:
        \[ 3^k \leq n^2 \]
        \item \textbf{Løsning med hensyn til $k$:} Vi tager logaritmen med grundtal $3$ på begge sider:
        \[ k \leq \log_3 n^2 \]
        \item \textbf{Reduktion:} Da $\log_3 n^2 = 2 \log_3 n$:
        \[ k \leq 2 \log_3 n \]
        \item \textbf{Antal iterationer:} Løkken udføres:
        \[ k = \lfloor 2 \log_3 n \rfloor + 1 \]
    \end{itemize}
    
    \textbf{Konklusion:} Løkken kører $\lfloor 2 \log_3 n \rfloor + 1$ gange, så algoritmens tidskompleksitet er $\boldsymbol{\mathcal{O}(\log n)}$.
    
\end{example}

\begin{example}

    Se på følgende algoritme:
    
    \begin{codebox}
        \Procname{$\proc{Loop5}(n)$}
        \li $i \gets 1$
        \li \While $i \leq n$ \Do
        \li     $j \gets 0$
        \li     \While $j \leq n$ \Do
        \li         $j \gets j + 1$
                \End
        \li     $i \gets 2 \times i$
            \End
    \end{codebox}
    
    Vi skal bestemme det samlede antal gange, den indre sætning $j \gets j + 1$ udføres.
    
    \begin{itemize}
        \item \textbf{Ydre løkke (\texttt{\While} $i \leq n$):}
        \begin{itemize}
            \item Variablen $i$ starter ved $1$ og fordobles hver gang: $i = 1, 2, 4, 8, \ldots$.
            \item Antallet af iterationer i den ydre løkke er:
            \[ k = \lfloor \log_2 n \rfloor + 1 \]
            \item Det skyldes, at løkken fortsætter, så længe $i \leq n$, og at $i$ vokser eksponentielt.
        \end{itemize}
        \item \textbf{Indre løkke (\texttt{\While} $j \leq n$):}
        \begin{itemize}
            \item For hver iteration af den ydre løkke initialiseres $j$ til $0$.
            \item Den indre løkke øger $j$ fra $0$ til $n$ og udføres $n + 1$ gange.
        \end{itemize}
        \item \textbf{Operationer i alt:}
        \begin{itemize}
            \item Den indre løkke kører $n + 1$ gange pr. iteration af den ydre løkke.
            \item Det samlede antal gange, $j \gets j + 1$ udføres:
            \[ (\lfloor \log_2 n \rfloor + 1) \times (n + 1) \]
            \item Ved reduktion får vi:
            \[ \mathcal{O}((\log n) \times n) = \mathcal{O}(n \log n) \]
        \end{itemize}
    \end{itemize}
    
    \textbf{Konklusion:} Algoritmen udfører $(\lfloor \log_2 n \rfloor + 1) \times (n + 1)$ operationer, så tidskompleksiteten er $\boldsymbol{\mathcal{O}(n \log n)}$.
    
\end{example}

\begin{example}

    Se på følgende algoritme:
    
    \begin{codebox}
        \Procname{$\proc{Loop6}(n)$}
        \li $i \gets 1$
        \li \While $i \leq n$ \Do
        \li     $j \gets i$
        \li     \While $j \leq n$ \Do
        \li         $j \gets j + 1$
                \End
        \li     $i \gets 2 \times i$
            \End
    \end{codebox}
    
    \textbf{Fuldstændig analyse af tidskompleksiteten:}
    
    Vi skal bestemme det samlede antal gange, den indre sætning $j \gets j + 1$ udføres.
    
    \begin{itemize}
        \item \textbf{Ydre løkke (\texttt{\While} $i \leq n$):}
        \begin{itemize}
            \item Variablen $i$ starter ved $1$ og fordobles hver gang: $i = 1, 2, 4, 8, \ldots$.
            \item Antallet af iterationer i den ydre løkke er:
            \[ k_{\text{max}} = \lfloor \log_2 n \rfloor + 1 \]
            \item Det skyldes, at løkken fortsætter, så længe $i \leq n$, og at $i$ vokser eksponentielt.
        \end{itemize}
        \item \textbf{Indre løkke (\texttt{\While} $j \leq n$):}
        \begin{itemize}
            \item For hver iteration af den ydre løkke initialiseres $j$ til den aktuelle værdi af $i$.
            \item Den indre løkke øger $j$ fra $i$ til $n$ og udføres $n - i + 1$ gange.
        \end{itemize}
        \item \textbf{Operationer i alt:}
        \begin{itemize}
            \item Det samlede antal gange, $j \gets j + 1$ udføres, er:
            \[
            \text{I alt} = \sum_{k=0}^{k_{\text{max}} - 1} (n - i_k + 1)
            \]
            hvor $i_k = 2^k$.
            \item Vi indsætter $i_k$:
            \[
            \text{I alt} = \sum_{k=0}^{\lfloor \log_2 n \rfloor} (n - 2^k + 1)
            \]
            \item Vi reducerer summen:
            \[
            \begin{aligned}
            \text{I alt} &= \sum_{k=0}^{\lfloor \log_2 n \rfloor} (n + 1 - 2^k) \\
            &= (\lfloor \log_2 n \rfloor + 1)(n + 1) - \sum_{k=0}^{\lfloor \log_2 n \rfloor} 2^k
            \end{aligned}
            \]
            \item Vi beregner summen af $2^k$:
            \[
            \sum_{k=0}^{\lfloor \log_2 n \rfloor} 2^k = 2^{\lfloor \log_2 n \rfloor + 1} - 1 \leq 2^{\log_2 n + 1} - 1 = 2n -1
            \]
            \item Derfor:
            \[
            \begin{aligned}
            \text{I alt} &= (\lfloor\log_2 n\rfloor+1)(n+1) - (2^{\lfloor\log_2 n\rfloor+1}-1) \\
            &\leq (\log_2 n+1)(n+1).
            \end{aligned}
            \]
            \item For store $n$ dominerer leddet $(\log_2 n)(n +1)$, så:
            \[
            \text{I alt} = \mathcal{O}(n \log n)
            \]
        \end{itemize}
    \end{itemize}
    
    \textbf{Konklusion:} Algoritmen udfører højst $(\log_2 n+1)(n+1)$ iterationer af den indre løkke, så tidskompleksiteten er $\boldsymbol{\mathcal{O}(n \log n)}$.

    \textbf{Forenklet analyse af tidskompleksiteten:}

    Vi kan forenkle analysen ved at undersøge den ydre og den indre løkkes opførsel hver for sig.

\begin{itemize}
    \item \textbf{Ydre løkke (\texttt{\While} $i \leq n$):}
    \begin{itemize}
        \item Variablen $i$ starter ved $1$ og fordobles hver gang:
        \[ i = 1, 2, 4, 8, \ldots, n \]
        \item Fordoblingen giver et logaritmisk antal iterationer:
        \[ \text{Antal iterationer i den ydre løkke} = \mathcal{O}(\log n) \]
    \end{itemize}
    \item \textbf{Indre løkke (\texttt{\While} $j \leq n$):}
    \begin{itemize}
        \item For hver værdi af $i$ starter $j$ ved $i$ og øges med $1$, indtil $j > n$.
        \item Antallet af iterationer i den indre løkke for et givet $i$ er:
        \[ n - i + 1 \leq n \]
        \item Det betyder, at den indre løkke kører cirka $n$ gange for hver værdi af $i$.
        \item Den indre løkke kører derfor $\mathcal{O}(n)$ gange for hver iteration af den ydre løkke.
    \end{itemize}
    \item \textbf{Samlet kompleksitet:}
    \begin{itemize}
        \item Den ydre løkke kører $\mathcal{O}(\log n)$ gange.
        \item For hver iteration af den ydre løkke kører den indre løkke $\mathcal{O}(n)$ gange.
        \item Det samlede antal operationer er derfor:
        \[ \mathcal{O}(\log n) \times \mathcal{O}(n) = \mathcal{O}(n \log n) \]
    \end{itemize}
\end{itemize}
    
\end{example}
    
\section{Analyse af kode}
Vi giver nu en række eksempler, der illustrerer, hvordan man analyserer kodestykker. Al koden er Java-kode, men principperne gælder for ethvert programmeringssprog.

\begin{example}
    
    \begin{lstlisting}[style=javaStyle, caption={Simpel funktion}, label={lst:java_simple_function}]
    public static int sumSquareDifference(int n) {
        int sum = (n * (n + 1)) / 2;
        int sumOfSquares = (n * (n + 1) * (2 * n + 1)) / 6;
        return (sum * sum) - sumOfSquares;
    }
    \end{lstlisting}

    Linje 2: \texttt{int sum = (n * (n + 1)) / 2;}
    \begin{itemize}
        \item \textbf{Tildelinger}: 
            \begin{itemize}
                \item Tildeling \texttt{sum} = ... : \textbf{1 tidsenhed}
            \end{itemize}
        \item \textbf{Regneoperationer}: 
            \begin{itemize}
                \item Multiplikation \texttt{n * (n + 1)}: \textbf{1 tidsenhed}
                \item Addition \texttt{n + 1}: \textbf{1 tidsenhed}
                \item Division \texttt{(n * (n + 1)) / 2}: \textbf{1 tidsenhed}
            \end{itemize}
        \item \textbf{I alt for linje 2}: \textbf{4 tidsenheder}
    \end{itemize}

    Linje 3: \texttt{int sumOfSquares = (n * (n + 1) * (2 * n + 1)) / 6;}
    \begin{itemize}
        \item \textbf{Tildelinger}:
            \begin{itemize}
                \item Tildeling \texttt{sumOfSquares} = ... : \textbf{1 tidsenhed}
            \end{itemize}
        \item \textbf{Regneoperationer}: 
            \begin{itemize}
                \item Multiplikation \texttt{n * (n + 1)}: \textbf{1 tidsenhed}
                \item Addition \texttt{n + 1}: \textbf{1 tidsenhed}
                \item Multiplikation \texttt{(n * (n + 1)) * (2 * n + 1)}: \textbf{1 tidsenhed}
                \item Multiplikation \texttt{2 * n}: \textbf{1 tidsenhed}
                \item Addition \texttt{2 * n + 1}: \textbf{1 tidsenhed}
                \item Division \texttt{(...) / 6}: \textbf{1 tidsenhed}
            \end{itemize}
        \item \textbf{I alt for linje 3}: \textbf{7 tidsenheder}
    \end{itemize}

    Linje 4: \texttt{return (sum * sum) - sumOfSquares;}
    \begin{itemize}
        \item \textbf{Regneoperationer}:
            \begin{itemize}
                \item Multiplikation \texttt{sum * sum}: \textbf{1 tidsenhed}
                \item Subtraktion \texttt{(sum * sum) - sumOfSquares}: \textbf{1 tidsenhed}
            \end{itemize}
        \item \textbf{Returnering} (skrivning til output eller returværdi): 
            \begin{itemize}
                \item Returværdi: \textbf{1 tidsenhed}
            \end{itemize}
        \item \textbf{I alt for linje 4}: \textbf{3 tidsenheder}
    \end{itemize}

    Opsummering af den samlede tid (T):
    \begin{itemize}
        \item \textbf{Linje 2}: 4 tidsenheder
        \item \textbf{Linje 3}: 7 tidsenheder
        \item \textbf{Linje 4}: 3 tidsenheder
    \end{itemize}

    Den samlede tid $T$ for at udføre funktionen \texttt{sumSquareDifference} er altså:

    \[
    T = 4 + 7 + 3 = 14 \text{ tidsenheder}
    \]

    Med denne detaljerede analyse kan vi beregne tiden $T$ præcist ved at tælle alle de enkelte operationer, og vi ser, at tidskompleksiteten er $\boldsymbol{\mathcal{O}(1)}$.

    Vi kunne også skrive analysen som kodekommentarer i funktionen for at dokumentere tidskompleksiteten af hver linje.

    \begin{lstlisting}[style=javaStyle, caption={Simpel funktion}, label={lst:java_simple_function2}]
        public static int sumSquareDifference(int n)
        {
            int sum = (n * (n + 1)) / 2;        // 4 time units
            int sumOfSquares = (n * (n + 1) * (2 * n + 1)) / 6;     // 7 time units
            return (sum * sum) - sumOfSquares;      // 3 time units
        } // Total time: T(n) = 4 + 7 + 3 = 14 time units so we get O(1)
    \end{lstlisting}
    
Vi foretrækker den sidste metode, fordi den giver en klar og kortfattet måde at dokumentere funktionens eller metodens tidskompleksitet på direkte i koden.
        
\end{example}

\begin{example}

I det følgende eksempel præsenterer vi en algoritme, der udfører iterativ binær søgning.

\begin{lstlisting}[style = javaStyle, caption={Iterativ binær søgning}, label={lst:java_binary_search_none}]
    public static int binarySearchIter(int[] a, int x) {
        int low = 0, high = a.length - 1;

        while (low <= high) {
            int mid = (low + high) / 2;
            if (x == a[mid]) {
                return mid;
            } 
            else if (x < a[mid]) {
                high = mid - 1;
            } 
            else {
                low = mid + 1;
            }
        }
        return -1;
    }
\end{lstlisting}

\begin{lstlisting}[style=javaStyle, caption={Iterativ binær søgning}, label={lst:java_binary_search}]
    public static int binarySearchIter(int[] a, int x) {
        int low = 0;
        int high = a.length - 1; // 2 time units (assignment: 1, subtraction: 1)

        // Loop runs log(n) times and makes 1 comparison per iteration
        while (low <= high) {
            int mid = (low + high) / 2; // For each iteration: 1 addition, 1 division, 1 assignment (3 time units)

            if (x == a[mid]) { // Comparison to check if element is found (1 time unit per iteration)
                return mid; // 1 time unit (return)
            } else if (x < a[mid]) { // Comparison to decide which half to search next (1 time unit per iteration)
                high = mid - 1; // For each iteration: 1 subtraction, 1 assignment (2 time units)
            } else { // If neither condition is satisfied
                low = mid + 1; // For each iteration: 1 addition, 1 assignment (2 time units)
            }
        }
        return -1; // 1 time unit (return)
    }
    // Detailed Time Complexity Analysis
    // ---------------------------------
    // Initialisations before the loop: 2 time units (assignments and subtraction)
    // For each iteration of the while loop (which runs log(n) times):
    // - Loop condition comparison: 1 time unit
    // - Midpoint calculation: 3 time units
    // - Comparisons to check element: 1 time unit
    // - One of the assignments for updating bounds (either low or high): 2 time units
    // Total time per iteration: 1 (while comparison) + 3 (mid calculation) + 1 (comparison) + 2 (update bounds) = 7 time units per iteration
    // The loop runs approximately log(n) times.
    // T(n) = 2 (initialisation) + 7 * log(n) (loop iterations) + 1 (final return)
    //      = 3 + 7 * log(n)
    // Ignoring constants and lower-order terms, we get: T(n) = O(log n)

\end{lstlisting}
    
\end{example}

Når vi analyserer kode, er vi oftest interesseret i tidskompleksiteten udtrykt med store $\mathcal{O}$-notation. Den viser, hvordan algoritmen skalerer med inputstørrelsen. Den detaljerede analyse af de tidsenheder, der bruges på hver kodelinje, er nyttig til at forstå algoritmens effektivitet og finde flaskehalse. Det kan være nyttigt, når man optimerer kode for at opnå bedre ydelse. Også i det tilfælde behøver vi kun at fokusere på det dominerende led i tidskompleksiteten, fordi det er det, der bestemmer algoritmens samlede vækstrate.

Vi analyserer desuden som regel, hvor tidsenhederne bruges i hver kodelinje. De følgende eksempler bruger denne fremgangsmåde.

\subsection*{Optimering af en algoritmes effektivitet}
Det sidste afsnit i dette kapitel -- og i bogen -- giver et eksempel på, hvordan man optimerer en algoritme. Vi bruger eksemplet med at finde summen af alle mulige par i et array til at illustrere, hvordan man kan forbedre en algoritmes effektivitet. Algoritmen optimeres fra en naiv fremgangsmåde til stadig mere effektive løsninger. Algoritmen gør følgende:

\begin{itemize}
    \item Algoritmen beregner \textbf{den samlede sum af alle mulige parvise summer} $\operatorname{arr}[i]+\operatorname{arr}[j]$ i arrayet \texttt{arr}.
    \item Summen løber over alle \emph{ordnede par} af indeks $(i,j)$. Både $(i,j)$ og $(j,i)$ medtages altså, og det samme gør par med $i=j$.
\end{itemize}

\begin{example} Her er en manuel illustration af summen af par:

    Antag, at $\operatorname{arr}=[1,2,3]$

   \textbf{ Manuel beregning af alle parsummer:}
    
   \begin{align*}
    \text{totalSum} = & \; (1+1)+(1+2)+(1+3)+ \\
                    & \; (2+1)+(2+2)+(2+3)+ \\
                    & \; (3+1)+(3+2)+(3+3) \\
                    = & \; 2+3+4+3+4+5+4+5+6 \\
                    = & \; 36
    \end{align*}
    
\end{example}

Lad os nu se på den \textbf{naive fremgangsmåde} til at løse problemet, hvor vi gennemløber alle mulige par i arrayet:

\begin{lstlisting}[style = javaStyle, caption={Naiv sum af par med analyse}, label={lst:java_naive_sum_pairs_analysis}]
    public static int naiveSumOfPairs(int[] arr) {
        int n = arr.length; // 1 time unit
        int totalSum = 0; // 1 time unit

        // Outer loop over 'i' from 0 to n - 1
        for (int i = 0; i < n; i++) { // 2n + 2
            // Inner loop over 'j' from 0 to n - 1
            for (int j = 0; j < n; j++) { // (2n + 2) * n
                totalSum += arr[i] + arr[j]; // 3n^2
            }
        }
        return totalSum; // 1 time unit
    }
    // Detailed Time Complexity Analysis
    // ---------------------------------
    // Initializations: 2 time units (n and totalSum assignments)
    // Outer loop: 2n + 2 time units
    // Inner loop 1: n * (2n + 2) time units
    // Inner loop body: 3n^2 time units
    // Return statement: 1 time unit
    //
    // Total Time: T(n) = 2 (initializations)
    //                    + 2n + 2
    //                    + 2n^2 + 2n
    //                    + 3n^2
    //                    + 1 (return)
    //                    = 5n^2 + 4n + 5
    // Dominating Term: 5n^2
    // Tight asymptotic complexity: Theta(n^2)
\end{lstlisting}

Algoritmen giver os en intuitiv løsning på problemet. Desværre er den meget ineffektiv for store input. Lad os nu se på en mere optimeret algoritme til samme problem. Analysen af den indre løkke bruger de summationsmetoder, der er samlet i Bilag \ref{app:appendix01}.

\begin{lstlisting}[style = javaStyle, caption={Optimeret sum af par med analyse}, label={lst:java_optimized_sum_pairs_analysis}]
    public static int optimizedSumOfPairs(int[] arr) {
        int n = arr.length;            // 1 time unit
        int totalSum = 0;              // 1 time unit
    
        // Outer loop over 'i' from 0 to n - 1
        for (int i = 0; i < n; i++) {  // 2n + 2 time units
            // Inner loop over 'j' from i to n - 1
            for (int j = i; j < n; j++) { // (2(n - i) + 2) per i, summed over i
                if (i == j) {              // Comparison: (n(n + 1)/2) time units
                    totalSum += arr[i] + arr[j]; // 3 time units, executed n times
                } else {
                    totalSum += 2 * (arr[i] + arr[j]); // 4 time units, executed (n^2 - n)/2 times
                }
            }
        }
        return totalSum; // 1 time unit
    }
    
    // Detailed Time Complexity Analysis
    // ---------------------------------
    // Initializations: 2 time units (n and totalSum assignments)
    // Outer loop: 2n + 2 time units (1 initialization, n + 1 comparisons, n increments)
    // Inner loop control: Sum of (2(n - i) + 2) for i from 0 to n - 1
    //     = Sum of (2n - 2i + 2) from i = 0 to n - 1
    //     = 2n^2 + 2n - n(n - 1)
    //     = n^2 + 3n time units
    // Inner loop body (if condition): 
    //     - Comparisons: (n(n + 1)/2) time units
    //     - if block (i == j): 3 time units, executed n times => 3n
    //     - else block (i != j): 4 time units, executed (n^2 - n)/2 times => 2n^2 - 2n
    //     Total inner loop body: n(n + 1)/2 + 3n + 2n^2 - 2n
    //     = (5n^2 + 3n)/2 time units
    // Return statement: 1 time unit
    //
    // Total Time: T(n) = 2 (initializations)
    //                    + 2n + 2 (outer loop)
    //                    + n^2 + 3n (inner loop control)
    //                    + (5n^2 + 3n)/2 (inner loop body)
    //                    + 1 (return)
    //                    = (7n^2 + 13n + 10)/2
    // Dominating Term: (7/2)n^2
    // Tight asymptotic complexity: Theta(n^2)
\end{lstlisting}

Denne algoritme er stadig $\Theta(n^2)$, men i vores enhedsomkostningsmodel er dens ledende led reduceret fra $5n^2$ til $\tfrac{7}{2}n^2$. Med de fuldstændige formler for antallet af operationer får vi ved $n=10^5$:

\begin{enumerate}
    \item $\proc{naiveSumOfPairs}$ (\(5n^2+4n+5\)):
    \[
    T(10^5)=5(10^5)^2+4(10^5)+5=50\,000\,400\,005
    \]

    \item $\proc{optimizedSumOfPairs}$ (\(\tfrac{7n^2+13n+10}{2}\)):
    \[
    T(10^5)=\frac{7(10^5)^2+13(10^5)+10}{2}=35\,000\,650\,005
    \]
\end{enumerate}

\[
\text{Forskel} = 50\,000\,400\,005-35\,000\,650\,005=14\,999\,750\,000
\]

For \(n=10^5\) udfører den naive version altså cirka \textbf{15 milliarder flere operationer} i denne model. Begge algoritmer har ikke desto mindre samme stramme asymptotiske kompleksitet, $\Theta(n^2)$.

Den næste algoritme er endnu mere effektiv med køretiden $\Theta(n)$.
\begin{lstlisting}[style = javaStyle, caption={Effektiv sum af par med analyse}, label={lst:java_efficient_sum_pairs_analysis}]
    public static int efficientSumOfPairs(int[] arr) {
        int n = arr.length;            // 1
        int totalSum = 0;              // 1
    
        for (int i = 0; i < n; i++) { //2n + 2
            totalSum += 2 * arr[i];   // 3n
        }
    
        return n * totalSum;           // 2
    }
        // Detailed Time Complexity Analysis
        // ---------------------------------
        // Initializations: 2 time units (n and totalSum assignments)
        // Loop: 2n + 2 time units (1 initialization, n + 1 comparisons, n increments)
        // Body: 3n time units (multiplication, addition and assignment, n iterations)
        // Return statement: 2 time units (multiplication and return)
        //
        // Total Time: T(n) = 2 (initializations)
        //                    + 2n + 2 (loop)
        //                    + 3n (body)
        //                    + 2 (return)
        //          = 5n + 6
        //
        // Tight asymptotic complexity: Theta(n)
\end{lstlisting}    

Denne algoritme er langt mere effektiv end de to foregående algoritmer. Dens køretid er $\Theta(n)$, hvilket asymptotisk er bedre end deres kvadratiske køretider. For inputstørrelsen \(n = 10^5\) er antallet af operationer:

\begin{enumerate}
    \item $\proc{efficientSumOfPairs}$ (\(5n+6\)):
    \[
    \text{Antal operationer} = 5n+6 = 5\cdot 10^5+6 = 500\,006
    \]
    \item $\proc{optimizedSumOfPairs}$ (\(\tfrac{7n^2+13n+10}{2}\)):
    \[
    \text{Antal operationer} = 35\,000\,650\,005
    \]
    \item $\text{Forskel} = 35\,000\,650\,005-500\,006=35\,000\,149\,999$
\end{enumerate}

Det illustrerer gevinsten ved at erstatte en $\Theta(n^2)$-algoritme med en $\Theta(n)$-algoritme. Den følgende algoritme har også køretiden $\Theta(n)$, men bruger færre operationer i kapitlets model:


\begin{lstlisting}[style = javaStyle, caption={Hurtig sum af par med analyse}, label={lst:java_fast_sum_pairs_analysis}]
    public static int fastSumOfPairs(int[] arr) {
        int n = arr.length;               // 1 time unit
        int sumOfArr = 0;                 // 1 time unit
    
        // Loop over 'i' from 0 to n - 1
        for (int i = 0; i < n; i++) {     // 2n + 2 time units (1 initialization, n + 1 comparisons, n increments)
            sumOfArr += arr[i];           // 2n time units (addition and assignment per iteration)
        }
    
        int totalSum = 2 * n * sumOfArr;  // 3 time units (2 multiplications, 1 assignment)
        return totalSum;                  // 1 time unit
    }
    // Updated Detailed Time Complexity Analysis
    // -----------------------------------------
    // Initializations: 2 time units (n and sumOfArr assignments)
    // Loop over 'i':
    //     - Loop control: 2n + 2 time units (1 initialization, n + 1 comparisons, n increments)
    //     - Loop body: 2n time units (sumOfArr += arr[i]; addition and assignment)
    // Total for loop: (2n + 2) + 2n = 4n + 2 time units
    // Computation of totalSum: 3 time units (2 multiplications, 1 assignment)
    // Return statement: 1 time unit
    //
    // Total Time: T(n) = 2 (initializations)
    //                    + 4n + 2 (loop over 'i')
    //                    + 3 (compute totalSum)
    //                    + 1 (return)
    //                    = 4n + 8
    // Dominating Term: 4n
    // Tight asymptotic complexity: Theta(n)
    
\end{lstlisting}  

I kapitlets model for optælling af operationer bruger $\proc{fastSumOfPairs}$ $4n+8$ operationer i stedet for $5n+6$. Ved $n=10^5$ er forskellen $99\,998$ operationer. Begge algoritmer har ikke desto mindre køretiden $\Theta(n)$, og den faktiske ydelse bør måles, hvis forskellen i den konstante faktor har betydning. Vi slutter med en sidste algoritme, der har samme asymptotiske kompleksitet og et lidt mindre antal operationer i denne model.

\begin{lstlisting}[style = javaStyle, caption={Ultraeffektiv sum af par med analyse}, label={lst:java_ultra_efficient_sum_pairs_analysis}]
    public static int ultraEfficientSumOfPairs(int[] arr) {
        int n = arr.length;               // 1 time unit
        int sumOfArr = 0;                 // 1 time unit
    
        // Loop over 'value' in 'arr'
        for (int value : arr) {           // 2n + 1 time units (n iterations)
            sumOfArr += value;            // 2n time units (addition and assignment per iteration)
        }
    
        int totalSum = 2 * n * sumOfArr;  // 3 time units (2 multiplications, 1 assignment)
        return totalSum;                  // 1 time unit
    }
    
    // Updated Detailed Time Complexity Analysis
    // -----------------------------------------
    // Initializations: 2 time units (n and sumOfArr assignments)
    // Loop over 'arr':
    //     - Loop control: 2n + 1 time units (1 initialization, n comparisons, n increments)
    //     - Loop body: 2n time units (sumOfArr += value; addition and assignment)
    // Total for loop: (2n + 1) + 2n = 4n + 1 time units
    // Computation of totalSum: 3 time units (2 multiplications, 1 assignment)
    // Return statement: 1 time unit
    //
    // Total Time: T(n) = 2 (initializations)
    //                    + (4n + 1) (loop over 'arr')
    //                    + 3 (compute totalSum)
    //                    + 1 (return)
    //                    = 4n + 7
    // Dominating Term: 4n
    // Tight asymptotic complexity: Theta(n)
    
\end{lstlisting}  


De to har identisk asymptotisk kompleksitet, og forskellen i konstanten er ubetydelig.

\textbf{Hvorfor $\proc{ultraEfficientSumOfPairs}$ er bedre:}
\begin{enumerate}
    \item \textbf{Læsbarhed:} Den udvidede \texttt{for-each}-løkke gennemløber elementerne direkte:
    \[
    \texttt{for (int value : arr) \{ sumOfArr += value; \}}
    \]
    Den er kortfattet, undgår indeksering og giver mindre risiko for fejl end en traditionel \texttt{for}-løkke:
    \[
    \texttt{for (int i = 0; i < n; i++) \{ sumOfArr += arr[i]; \}}
    \]
    \item \textbf{Vedligeholdelse:} Udvidede løkker er lettere at ændre og tilpasse til forskellige datastrukturer.
    \item \textbf{Moderne praksis:} Den følger nutidens Java-standarder og gør koden og hensigten tydeligere.
\end{enumerate}

Selv om de to funktioner yder omtrent lige godt, foretrækkes $\proc{ultraEfficientSumOfPairs}$ på grund af sin læsbarhed, enkelhed og overensstemmelse med moderne programmeringspraksis, hvilket gør den til en mere elegant løsning.

